\chapter{Buenas Prácticas de Desarrollo Backend}
\label{ch:buenas-practicas}

Este capítulo aborda el indicador de mayor ponderación del informe (16\%): el desarrollo
de componentes de backend aplicando buenas prácticas de diseño y arquitectura, el uso de
Spring y Maven, y el trabajo colaborativo mediante Git.

\section{Patrón Controller-Service-Repository}

Los diez microservicios de BookPoint aplican de forma consistente el patrón
\textbf{Controller-Service-Repository (CSR)}, separando estrictamente tres
responsabilidades:

\begin{itemize}
    \item \textbf{Controller}: recibe la petición HTTP, delega en el service y traduce
    el resultado a una respuesta REST (código de estado, cuerpo, enlaces HATEOAS). No
    contiene lógica de negocio.
    \item \textbf{Service}: concentra la lógica de negocio, las validaciones cruzadas
    contra otros microservicios (vía Feign) y las reglas del dominio (por ejemplo,
    verificar stock antes de confirmar un pedido).
    \item \textbf{Repository}: acceso a datos mediante Spring Data JPA
    (\texttt{JpaRepository}), sin lógica de negocio.
\end{itemize}

El siguiente ejemplo, tomado del microservicio de catálogo, muestra la capa de servicio
completa de \texttt{LibroServiceImpl}: no contiene código HTTP ni SQL, solo la
traducción entre el DTO de entrada y la entidad persistida, y el manejo del caso "no
encontrado" mediante una excepción de negocio propia:

\begin{lstlisting}[language=Java, caption={LibroServiceImpl.java -- capa de servicio de catálogo}]
@Service
@RequiredArgsConstructor
public class LibroServiceImpl implements LibroService {

    private final LibroRepository libroRepository;

    @Override
    public Libro guardarLibro(LibroDTO libroDTO) {
        Libro libro = new Libro();
        libro.setTitulo(libroDTO.getTitulo());
        libro.setAutor(libroDTO.getAutor());
        libro.setPrecio(libroDTO.getPrecio());
        return libroRepository.save(libro);
    }

    @Override
    public Libro obtenerPorId(Long id) {
        return libroRepository.findById(id)
            .orElseThrow(() -> new RecursoNoEncontradoException(
                "Libro no encontrado con id: " + id));
    }
}
\end{lstlisting}

Esta separación se replica de forma idéntica en los diez servicios, lo que permite que
cualquier persona familiarizada con uno de ellos pueda orientarse rápidamente en los
demás -- una consecuencia directa de mantener una convención uniforme en todo el
proyecto.

\section{Modularidad, mantenibilidad y escalabilidad}

\subsection{Base de datos por servicio}

Cada microservicio posee su propia base de datos MySQL, sin compartir esquemas ni
tablas con otros servicios (Tabla~\ref{tab:bd-servicio}). Esto evita el acoplamiento de
datos característico de los monolitos -- como el descrito en el caso BookPoint Chile,
donde ventas, inventario, despacho y atención de clientes conviven en un mismo sistema
-- y permite que cada equipo pueda evolucionar el modelo de datos de su dominio sin
coordinar migraciones con el resto.

\begin{table}[h]
\centering
\begin{tabular}{ll}
\toprule
\textbf{Microservicio} & \textbf{Base de datos} \\
\midrule
catálogo & bookpoint\_catalogo\_db \\
inventario & bookpoint\_inventario\_db \\
pedidos & bookpoint\_pedidos\_db \\
clientes & bookpoint\_clientes\_db \\
sucursales & bookpoint\_sucursales\_db \\
carrito & bookpoint\_carrito\_db \\
pagos & bookpoint\_pagos\_db \\
envíos & bookpoint\_envios\_db \\
notificaciones & bookpoint\_notificaciones\_db \\
reseñas & bookpoint\_resenas\_db \\
\bottomrule
\end{tabular}
\caption{Una base de datos independiente por microservicio}
\label{tab:bd-servicio}
\end{table}

\subsection{Separación DTO / entidad}

En los servicios donde la entrada requiere validación estricta (por ejemplo,
\texttt{catálogo}, \texttt{inventario} y \texttt{pedidos}), se utiliza un DTO anotado
con \textit{Bean Validation} (\texttt{@NotBlank}, \texttt{@NotNull}, \texttt{@Positive})
en lugar de exponer la entidad JPA directamente en el controlador. Esto evita que un
cliente de la API pueda, por ejemplo, enviar un \texttt{id} propio al crear un recurso
y sobrescribir accidentalmente uno existente -- un error real detectado y corregido
durante el desarrollo (ver sección~\ref{sec:evolucion}).

\subsection{Manejo centralizado de excepciones}

Cada microservicio expone un \texttt{@RestControllerAdvice} propio
(\texttt{GlobalExceptionHandler}) que traduce las excepciones de negocio en respuestas
HTTP coherentes, evitando repetir bloques \texttt{try/catch} en cada controlador:

\begin{lstlisting}[language=Java, caption={GlobalExceptionHandler.java (catálogo) -- fragmento}]
@RestControllerAdvice
public class GlobalExceptionHandler {

    @ExceptionHandler(RecursoNoEncontradoException.class)
    public ResponseEntity<ErrorResponseDTO> handleRecursoNoEncontrado(
            RecursoNoEncontradoException ex) {
        ErrorResponseDTO errorDTO = new ErrorResponseDTO(
                ex.getMessage(), HttpStatus.NOT_FOUND.value(),
                LocalDateTime.now(), null);
        return new ResponseEntity<>(errorDTO, HttpStatus.NOT_FOUND);
    }

    @ExceptionHandler(IllegalArgumentException.class)
    public ResponseEntity<ErrorResponseDTO> handleArgumentoInvalido(
            IllegalArgumentException ex) {
        ErrorResponseDTO errorDTO = new ErrorResponseDTO(
                ex.getMessage(), HttpStatus.BAD_REQUEST.value(),
                LocalDateTime.now(), null);
        return new ResponseEntity<>(errorDTO, HttpStatus.BAD_REQUEST);
    }
}
\end{lstlisting}

Este patrón asegura que un recurso inexistente responda siempre con \texttt{404}, una
violación de regla de negocio con \texttt{400}, y un error inesperado con \texttt{500}
-- de forma consistente en los diez servicios, en vez de depender de que cada
desarrollador recuerde hacerlo caso a caso.

\section{Spring Boot y Maven}

El proyecto está organizado como un \textbf{POM multi-módulo} de Maven: un
\texttt{pom.xml} raíz de tipo \texttt{pom} declara los doce módulos (diez
microservicios de negocio, \texttt{eureka-server} y \texttt{gateway}) y centraliza en
\texttt{dependencyManagement} las versiones compartidas -- \textbf{Spring Boot 4.1},
\textbf{Java 21} y \textbf{Spring Cloud 2025.1.2} -- de modo que ningún módulo hijo
necesita fijar versiones por su cuenta:

\begin{lstlisting}[language=XML, caption={pom.xml raíz -- fragmento relevante}]
<properties>
    <java.version>21</java.version>
    <spring-cloud.version>2025.1.2</spring-cloud.version>
</properties>

<dependencyManagement>
    <dependencies>
        <dependency>
            <groupId>org.springframework.cloud</groupId>
            <artifactId>spring-cloud-dependencies</artifactId>
            <version>${spring-cloud.version}</version>
            <type>pom</type>
            <scope>import</scope>
        </dependency>
    </dependencies>
</dependencyManagement>
\end{lstlisting}

Además de \texttt{spring-boot-starter-web} y \texttt{spring-boot-starter-data-jpa}
(presentes en los diez servicios), el proyecto utiliza:

\begin{itemize}
    \item \textbf{Spring Cloud OpenFeign}: comunicación declarativa entre
    microservicios (capítulo~\ref{ch:rest}).
    \item \textbf{Spring Cloud Netflix Eureka Client}: registro y descubrimiento de
    servicios.
    \item \textbf{Spring HATEOAS}: navegabilidad de la API (capítulo~\ref{ch:hateoas}).
    \item \textbf{springdoc-openapi}: documentación Swagger/OpenAPI por servicio.
    \item \textbf{JaCoCo Maven Plugin}: medición de cobertura de pruebas
    (capítulo~\ref{ch:testing}).
    \item \textbf{Datafaker}: generación de datos de prueba realistas al levantar cada
    servicio (ver sección~\ref{sec:evolucion}).
\end{itemize}

\section{Trabajo colaborativo con Git}

Este proyecto fue desarrollado íntegramente en solitario: aunque la pauta permite
equipos de hasta tres integrantes, todo el desarrollo backend -- diseño, código,
pruebas y correcciones -- fue trabajo individual, sin otros integrantes aportando
código al repositorio.

El trabajo se hizo directo sobre la rama \texttt{main}, sin ramas separadas por
funcionalidad. El desarrollo comenzó por los servicios núcleo del flujo comercial
(catálogo, inventario y pedidos), construidos hasta tener la comunicación Feign básica
funcionando entre ellos; a partir de ahí se fue extendiendo el resto de los
microservicios y, en tareas más mecánicas y repetitivas -- aplicar el mismo cambio de
forma consistente a lo largo de los diez servicios, como la incorporación de HATEOAS o
la construcción de la colección Postman --, se aceleró el trabajo con herramientas de
asistencia de desarrollo.

El repositorio del proyecto utiliza Git como sistema de control de versiones, con una
convención de mensajes de commit en español, explicando siempre el \textit{porqué} del
cambio y no solo el \textit{qué}. Los cambios no relacionados se separan en commits
distintos en lugar de agruparse en una sola entrega masiva -- por ejemplo, la sesión de
trabajo dedicada a corregir el mapeo de errores HTTP, agregar HATEOAS y sumar
validaciones cruzadas vía Feign se dividió en tres commits independientes, cada uno
enfocado en un solo cambio lógico:

\begin{lstlisting}[caption={Extracto de \texttt{git log --oneline}}]
47bacbb feat: implementar HATEOAS en los 9 microservicios que faltaban
9c05bdb feat: agregar validaciones cruzadas via Feign en carrito,
        resenas, pagos y envios
ccaf47d fix: mapear "no encontrado" a HTTP 404 en los 10 microservicios
\end{lstlisting}

\section{Evolución del proyecto: correcciones realizadas durante el semestre}
\label{sec:evolucion}

A diferencia de las evaluaciones parciales anteriores -- que consistieron
exclusivamente en entrega de código y defensa técnica en vivo, sin un informe escrito
--, este documento sí exige explicitar las correcciones realizadas a lo largo del
semestre. El historial de Git del proyecto permite reconstruir esa evolución con
precisión, sin depender de la memoria:

\begin{description}
    \item[Mapeo de errores HTTP.] Hasta antes de la corrección, cualquier recurso no
    encontrado (por ejemplo, un libro con un \texttt{id} inexistente) devolvía
    \texttt{500 Internal Server Error} en lugar de \texttt{404}, porque no existía un
    manejo de excepciones dedicado. Se creó una excepción
    \texttt{RecursoNoEncontradoException} y un \texttt{GlobalExceptionHandler} por
    servicio (reutilizando el que ya existía en catálogo en vez de duplicarlo), lo que
    también reactivó bloques \texttt{catch (FeignException.NotFound)} que hasta
    entonces eran código muerto.
    \item[HATEOAS.] Inicialmente solo el microservicio de catálogo implementaba
    HATEOAS. Se extendió el mismo patrón (\texttt{EntityModel}/\texttt{CollectionModel}
    con enlaces \textit{self}) a los nueve microservicios restantes.
    \item[Validaciones cruzadas vía Feign.] Se agregaron validaciones adicionales entre
    servicios que aún no las tenían: \texttt{carrito} valida contra \texttt{clientes} y
    \texttt{catálogo}; \texttt{reseñas} valida contra \texttt{catálogo} y
    \texttt{clientes}; \texttt{pagos} y \texttt{envíos} validan contra \texttt{pedidos}.
    \item[Cobertura de pruebas.] Se incorporó JaCoCo al POM raíz (antes no existía
    medición real de cobertura) y se agregó una clase de prueba
    \texttt{@WebMvcTest} por servicio, que verifica los códigos HTTP reales devueltos
    por cada endpoint. La cobertura de la lógica de negocio (controladores y
    servicios) es hoy de \textbf{93,6\%} (capítulo~\ref{ch:testing}).
    \item[Relación JPA real entre entidades.] Se agregó al microservicio de envíos una
    segunda entidad, \texttt{HistorialEstado}, en una relación real
    \texttt{@OneToMany}/\texttt{@ManyToOne} con \texttt{Envio}, que registra cada
    cambio de estado del envío. Se optó por este servicio en particular porque la
    relación es puramente aditiva: no modifica el contrato de ningún endpoint
    existente.
    \item[Datos de prueba con Datafaker.] Se incorporó un \texttt{DataSeeder} por
    servicio (\texttt{CommandLineRunner}) que genera datos de ejemplo realistas la
    primera vez que cada tabla está vacía, reutilizando la propia capa de servicio (y
    por lo tanto las validaciones Feign existentes) en lugar de insertar filas
    directamente en el repositorio.
    \item[Validación con Postman.] Se construyó una colección Postman que cubre CRUD y
    comunicación REST entre microservicios para los diez servicios, verificada en vivo
    contra el stack completo levantado en Docker antes de darla por válida
    (capítulos~\ref{ch:crud} y~\ref{ch:rest}).
    \item[Autenticación en clientes.] Se agregó JWT + BCrypt al microservicio de
    clientes, protegiendo únicamente las operaciones de escritura sobre el propio
    perfil (sección~\ref{ch:etica}). Se optó por este alcance acotado en vez de
    proteger todo el microservicio, precisamente porque \texttt{carrito} y
    \texttt{reseñas} dependen de poder leerlo vía Feign sin un token de usuario.
    \item[Reportes de ventas.] Se agregaron reportes de ventas por sucursal y por
    autor en \texttt{pedidos} (sección~\ref{ch:crud}), cerrando parcialmente la
    exclusión de reportes documentada originalmente en la sección~\ref{ch:arquitectura}.
\end{description}
